\documentclass[UTF8,a4paper,11pt,fontset=windows]{ctexart}
\usepackage[margin=21mm,top=22mm,bottom=22mm,headheight=15pt]{geometry}
\usepackage{xcolor}
\usepackage{fontspec,graphicx,amsmath,booktabs,array,longtable,enumitem,fancyhdr,tikz,tcolorbox,hyperref}
\usetikzlibrary{arrows.meta,positioning,fit,backgrounds,calc}
\definecolor{Navy}{HTML}{17324D}
\definecolor{Teal}{HTML}{007F82}
\definecolor{Gold}{HTML}{B77623}
\definecolor{Ink}{HTML}{253445}
\definecolor{Muted}{HTML}{607386}
\definecolor{Pale}{HTML}{F0F6F8}
\definecolor{Rule}{HTML}{CDDCE4}
\hypersetup{colorlinks=true,linkcolor=Navy,urlcolor=Teal,pdftitle={业务决策节点轻量化与双闭环 RSI 自优化方案},pdfauthor={},bookmarksnumbered=true}
\ctexset{section={format=\Large\bfseries\color{Navy},beforeskip=15pt,afterskip=10pt},subsection={format=\large\bfseries\color{Teal},beforeskip=12pt,afterskip=7pt}}
\setlength{\parindent}{2em}
\setlength{\parskip}{4pt}
\linespread{1.18}
\setlist{nosep,leftmargin=2em,topsep=5pt,itemsep=3pt}
\setlength{\emergencystretch}{3em}
\pagestyle{fancy}\fancyhf{}
\fancyhead[L]{\small\color{Muted}业务决策节点轻量化}
\fancyhead[R]{\small\color{Muted}双闭环 RSI 自优化方案}
\fancyfoot[L]{\scriptsize\color{Muted}参赛申请优化版 · 技术方案}
\fancyfoot[R]{\small\color{Navy}\thepage}
\renewcommand{\headrulewidth}{0.3pt}
\renewcommand{\footrulewidth}{0pt}
\newtcolorbox{insight}[1]{colback=Pale,colframe=Rule,boxrule=.5pt,arc=2mm,left=3mm,right=3mm,top=2mm,bottom=2mm,title={#1},colbacktitle=Navy,coltitle=white,fonttitle=\bfseries}
\newcommand{\figcap}[1]{\refstepcounter{figure}\par\smallskip{\centering\small\color{Muted}图\thefigure\quad #1\par}\medskip}
\tikzset{box/.style={draw=Rule,fill=Pale,rounded corners=2mm,align=center,text width=34mm,minimum height=12mm,font=\small,inner sep=4pt},accent/.style={box,draw=Teal,fill=Teal!8},dark/.style={box,fill=Navy,draw=Navy,text=white},arr/.style={-{Stealth[length=2mm]},thick,draw=Muted},lab/.style={font=\scriptsize,text=Muted,fill=white,inner sep=1pt}}
\begin{document}
\color{Ink}
\begin{titlepage}
\thispagestyle{empty}
\begin{tikzpicture}[remember picture,overlay]
\fill[Navy] (current page.north west) rectangle ([yshift=-62mm]current page.north east);
\fill[Teal] ([yshift=-62mm]current page.north west) rectangle ([yshift=-65mm]current page.north east);
\node[anchor=north west,text=white,font=\small] at ([xshift=22mm,yshift=-21mm]current page.north west) {技术方案 / 参赛申请优化版};
\node[anchor=north west,text=white,font=\fontsize{27}{36}\selectfont\bfseries,align=left] at ([xshift=22mm,yshift=-34mm]current page.north west) {让业务判断更轻，让优化持续发生};
\end{tikzpicture}\par
\vspace*{58mm}
{\fontsize{25}{36}\selectfont\bfseries\color{Navy}业务决策节点轻量化\par}
\vspace{4mm}
{\fontsize{23}{34}\selectfont\bfseries\color{Navy}与双闭环 RSI 自优化方案\par}
\vspace{7mm}
{\large\color{Muted}以业务质量为约束，以完整任务收益为依据\par}
\vspace{13mm}
\begin{tikzpicture}
\node[dark,text width=40mm,minimum height=23mm] (a) {节点轻量化\\[-1pt]\footnotesize 分类与选项打分};
\node[accent,text width=40mm,minimum height=23mm,right=8mm of a] (b) {提示优化闭环\\[-1pt]\footnotesize 固定权重，改进描述};
\node[box,text width=40mm,minimum height=23mm,right=8mm of b] (c) {业务微调闭环\\[-1pt]\footnotesize 冻结任务，补强能力};
\draw[arr] (a)--(b);\draw[arr] (b)--(c);
\end{tikzpicture}
\vspace{11mm}
\begin{insight}{方案主张}
保留必要的内容生成，将边界清晰的判断任务交给合适的决策后端；通过真实错误驱动的两条离线闭环，形成可评测、可追溯、可回退的持续优化机制。
\end{insight}
\vfill
{\small\color{Muted}文档状态：方案设计与验证规划，尚未提供实测成果。}
\end{titlepage}
\pagenumbering{roman}
\section*{方案速览}
\addcontentsline{toc}{section}{方案速览}
\begin{insight}{一句话理解}
用更适合判断任务的推理方式承接业务节点，在质量达标的前提下减少生成开销；再以提示优化和可选微调持续提高可靠承接能力。
\end{insight}
\begin{description}[style=nextline,leftmargin=2em,labelwidth=2em]
\item[为什么做] 智能助理中的意图分类、改写判断和证据筛选，往往只需输出类别、是非或分数，却可能消耗一次完整的生成式调用。
\item[从哪里开始] 首先验证意图识别，重点区分消费者退款、商户结算和支付状态；问题改写与检索片段筛选作为扩展场景。
\item[如何持续改进] 先诊断输入、标签与策略问题；描述不清时优化提示，领域能力不足且训练条件具备时再比较微调。两条闭环均在线下运行。
\item[怎样判断有效] 先满足预先约定的质量约束，再比较完整链路时延与每成功任务总成本。新增决策、回退、运维和离线优化摊销均计入成本。
\end{description}

\clearpage
\tableofcontents
\clearpage
\pagenumbering{arabic}


\section{背景与目标}

智能助理的一条业务链路，往往需要先识别用户意图，再判断是否改写问题、检索材料是否有用，最后生成回答。其中不少节点只需输出一个类别、是非结果或分数，却仍调用通用生成式大模型。


例如，用户问“退款已经申请了，怎么还没到账？”，意图识别节点只需要返回“退款进度查询”。当这类判断频繁出现或多个节点串联时，推理费用与等待时间会不断累积。与此同时，类别描述和提示词主要靠人工调整，业务错误尚未形成统一的优化反馈。


本方案将适合的判断节点交给决策模型处理，并建立提示优化、业务微调两条离线闭环。目标是在业务质量达标的前提下，降低完整任务的总成本与时延，同时形成可复用的节点优化能力。


项目围绕“用决策推理替代生成式判断、减少不必要的生成调用、减少回答模型的冗余输入”三条路径开展优化；通过双闭环 RSI 持续改善决策模型的业务适配能力，形成可对照、可验证、可回退的降本提速方案。


\section{具体使用场景}

\subsection{主场景：智能助理意图识别}

用户问题进入智能助理后，系统需要识别其业务诉求，再进入退款查询、商户结算查询等对应流程。决策模型承接的是路由前的判断，不负责查询业务数据，也不直接执行业务操作。


{\small\setlength{\tabcolsep}{4pt}\renewcommand{\arraystretch}{1.28}
\begin{longtable}{>{\raggedright\arraybackslash}p{0.220\dimexpr\linewidth-6\tabcolsep\relax}>{\raggedright\arraybackslash}p{0.380\dimexpr\linewidth-6\tabcolsep\relax}>{\raggedright\arraybackslash}p{0.400\dimexpr\linewidth-6\tabcolsep\relax}}
\caption{用户意图识别示例及下游处理}\label{tab:intent-examples}\\
\toprule
\textbf{用户表达示例} & \textbf{期望识别结果} & \textbf{后续处理} \\
\midrule
\endfirsthead
\toprule
\textbf{用户表达示例} & \textbf{期望识别结果} & \textbf{后续处理} \\
\midrule
\endhead
\midrule\multicolumn{3}{r}{\scriptsize\color{Muted}续下页}\\\endfoot
\bottomrule\endlastfoot
“退款已经申请了，怎么还没到账？” & 退款进度查询 & 进入退款查询流程 \\
“今天的收款为什么还没结算？” & 商户结算查询 & 进入结算查询流程 \\
“这笔订单到底付成功没有？” & 支付状态查询 & 进入支付状态查询流程 \\
前文说“我申请了退款”，当前问“什么时候到？” & 结合前文识别退款进度查询 & 使用必要的对话上下文 \\
没有前文，只问“什么时候到？” & 信息不足，需要澄清 & 询问所指业务 \\
问题明确，但不属于支持范围 & 未知意图 & 转原有兜底流程 \\
\end{longtable}}

“退款到账”与“结算到账”都有到账概念，但业务对象不同，是需要重点验证的混淆边界。多意图输入也应明确采用多标签、逐项判断或澄清，不能强行套用单选分类。


输入：当前用户问题、分类必需的对话上下文、经授权的结构化信息。


输出：类别 ID、各选项原始概率/分数、接受/回退状态、版本与耗时。




提示优化应更清晰表达这种差异，而非把任一类别扩展成"所有到账问题"来提高局部分数。


\textbf{替换前后链路：}



\begin{center}\resizebox{\linewidth}{!}{\begin{tikzpicture}
\node[dark,text width=53mm] (u) {用户问题与必要上下文};
\node[box,below left=10mm and 8mm of u,text width=48mm] (old) {现有路径\\生成式大模型判断};
\node[accent,below right=10mm and 8mm of u,text width=48mm] (new) {候选路径\\决策模型判断与接受策略};
\node[accent,below=17mm of new,text width=24mm] (ask) {信息不足\\澄清问题};
\node[box,left=5mm of ask,text width=28mm] (fb) {不可靠或异常\\原模型／既有兜底};
\node[dark,below left=66mm and 8mm of u,text width=28mm] (r) {业务路由};
\draw[arr] (u)-|(old);\draw[arr] (u)-|(new);
\draw[arr] (old.south) -- ++(-8mm,0) |- (r.west);
\draw[arr] (new)--node[lab,right]{缺少信息}(ask);
\draw[arr] (new.south west)--node[lab,sloped,above]{不满足}(fb.north);
\draw[arr] (new.east)--++(9mm,0)|-node[lab,near end,above]{满足接受条件}(r.east);
\draw[arr] (fb.west) -| (r.north east);
\end{tikzpicture}}\end{center}
\figcap{意图识别替换路径：接受、澄清与回退分别处理}


图中两条路径用于比较现有与候选方案；正式运行时按批准配置选择路径。候选方案需验证完整路由结果，不能只展示模型返回的类别。


\textbf{场景应用举例（原本如何 → 用本方案预期如何）}


以"退款已经申请了，怎么还没到账？"为例：


{\small\setlength{\tabcolsep}{4pt}\renewcommand{\arraystretch}{1.28}
\begin{longtable}{>{\raggedright\arraybackslash}p{0.220\dimexpr\linewidth-6\tabcolsep\relax}>{\raggedright\arraybackslash}p{0.380\dimexpr\linewidth-6\tabcolsep\relax}>{\raggedright\arraybackslash}p{0.400\dimexpr\linewidth-6\tabcolsep\relax}}
\caption{意图识别流程优化前后对照}\label{tab:intent-flow}\\
\toprule
\textbf{环节} & \textbf{原本流程} & \textbf{用本方案后（预期）} \\
\midrule
\endfirsthead
\toprule
\textbf{环节} & \textbf{原本流程} & \textbf{用本方案后（预期）} \\
\midrule
\endhead
\midrule\multicolumn{3}{r}{\scriptsize\color{Muted}续下页}\\\endfoot
\bottomrule\endlastfoot
判断方式 & 调用通用生成式大模型，让模型输出类别名称 & 调用轻量决策模型，直接输出类别 ID 与各选项概率 \\
耗时 & 生成式推理串联，等待首 token 与完整文本 & 采用支持分类或选项打分的后端，预期减少生成式解码开销，实际时延待测 \\
成本 & 生成式调用按实际输入、输出及计费规则核算，高频调用累积成本高 & 按所选后端的实际计费或部署成本核算，验证单次决策费用是否下降 \\
易混淆处理 & 依赖人工调整提示词区分"退款到账"与"结算到账"，调整滞后 & 离线提示优化迭代类别描述，经评测与审批后发布优化版本 \\
未知/澄清 & 依赖现有接受与兜底策略，实际处理方式需先核查 & 置信不足时进入 CLARIFY 或 FALLBACK，回到原流程，不强行归类 \\
\end{longtable}}

\subsection{扩展场景：是否需要问题改写}

在多轮问答中，用户说“那什么时候到？”，系统可能需要结合前文补全为“已申请的退款预计什么时候到账？”，再进行检索。


决策模型只判断是否进入改写路径；需要改写时，仍由生成模型产出新问题。对于已经完整明确的问题，可直接检索，减少不必要的生成调用。


该场景的价值取决于检索和回答是否改善。漏掉必要改写可能损害召回，多做改写会增加成本，因此不能单纯追求降低改写率。


评测同时比较原始与改写后检索的关键证据和最终回答质量，依据漏改写的质量损失与多改写的额外费用设定策略。


\textbf{场景应用举例}


以多轮对话中"那笔呢？"（指代前文某笔交易）为例：


{\small\setlength{\tabcolsep}{4pt}\renewcommand{\arraystretch}{1.28}
\begin{longtable}{>{\raggedright\arraybackslash}p{0.220\dimexpr\linewidth-6\tabcolsep\relax}>{\raggedright\arraybackslash}p{0.380\dimexpr\linewidth-6\tabcolsep\relax}>{\raggedright\arraybackslash}p{0.400\dimexpr\linewidth-6\tabcolsep\relax}}
\caption{问题改写场景优化前后对照}\label{tab:rewrite-flow}\\
\toprule
\textbf{环节} & \textbf{原本流程} & \textbf{用本方案后（预期）} \\
\midrule
\endfirsthead
\toprule
\textbf{环节} & \textbf{原本流程} & \textbf{用本方案后（预期）} \\
\midrule
\endhead
\midrule\multicolumn{3}{r}{\scriptsize\color{Muted}续下页}\\\endfoot
\bottomrule\endlastfoot
改写触发 & 在默认改写的流程中，每条问题都送生成式大模型改写，再检索 & 决策模型先判断是否需要改写，确认无需改写的完整问题直接检索 \\
成本 & 每次调用都付出改写生成费用，包括无需改写的简单问题 & 仅需改写的问题才调用大模型，预期改写调用占比下降、总成本降低 \\
风险 & 若缺少统一评测，漏改写与多改写主要依靠人工经验权衡 & 用"必要改写召回 / 多余改写率"量化，按质量损失与额外费用设定阈值 \\
回退 & 需核查现有流程是否具备检索失败后的补救机制 & 判错可经反馈回流进入错误分析池，驱动提示或策略迭代 \\
\end{longtable}}

\subsection{扩展场景：RAG 检索片段筛选}

检索返回多段退款规则、商户结算说明和产品材料后，决策模型判断哪些片段对当前问题有帮助，再交给回答模型。


沿用已有分块与检索，先做"是否对回答有帮助"。


输入应包含问题、片段及必要来源和上下文；关键证据、例外条款和冲突材料需保守保留。主题相关不等于能支撑答案，更不等于答案正确。


此场景应与已有 reranker 对照。如果专用 reranker 在同等条件下更好、更便宜，就保留它。项目不要求所有节点都使用同一类模型。


\textbf{场景应用举例}


以用户问"退款多久到账"，检索返回 5 个候选片段为例：


{\small\setlength{\tabcolsep}{4pt}\renewcommand{\arraystretch}{1.28}
\begin{longtable}{>{\raggedright\arraybackslash}p{0.220\dimexpr\linewidth-6\tabcolsep\relax}>{\raggedright\arraybackslash}p{0.380\dimexpr\linewidth-6\tabcolsep\relax}>{\raggedright\arraybackslash}p{0.400\dimexpr\linewidth-6\tabcolsep\relax}}
\caption{检索片段筛选场景优化前后对照}\label{tab:retrieval-flow}\\
\toprule
\textbf{环节} & \textbf{原本流程} & \textbf{用本方案后（预期）} \\
\midrule
\endfirsthead
\toprule
\textbf{环节} & \textbf{原本流程} & \textbf{用本方案后（预期）} \\
\midrule
\endhead
\midrule\multicolumn{3}{r}{\scriptsize\color{Muted}续下页}\\\endfoot
\bottomrule\endlastfoot
片段筛选 & 5 个片段全部送入生成式大模型上下文，由模型自行取舍 & 决策模型先标注"是否对回答有帮助"，仅保留关键证据片段 \\
上下文成本 & 冗余片段增加输入 token，推高生成费用与时延 & 裁剪后输入缩短，预期减少回答模型输入费用；完整链路时延需扣除筛选开销后验证 \\
关键证据 & 例外条款、冲突信息可能被模型忽略 & 保守保留例外与冲突片段，减少漏答风险 \\
对照 & 若缺少独立评测，则难以判断裁剪是否损害质量 & 用"证据召回 / 误删 / 保留比例"量化，与现有 reranker 公平对照后择优 \\
\end{longtable}}

\section{降本提速机制与预期成效}


\begin{center}\begin{tikzpicture}
\node[accent,text width=43mm,minimum height=22mm] (a) {01\quad 替换判断\\\footnotesize 分类／选项打分\\替代生成式判断};
\node[accent,right=6mm of a,text width=43mm,minimum height=22mm] (b) {02\quad 减少调用\\\footnotesize 先判断改写必要性\\跳过多余生成};
\node[accent,right=6mm of b,text width=43mm,minimum height=22mm] (c) {03\quad 缩短输入\\\footnotesize 筛选支撑性证据\\减少冗余上下文};
\node[dark,below=8mm of b,text width=147mm] (q) {共同前提：质量达标，且节省覆盖新增判断、筛选、回退与部署开销};
\draw[arr](a)--(q);\draw[arr](b)--(q);\draw[arr](c)--(q);
\end{tikzpicture}\end{center}
\figcap{三条收益路径可以独立采用，净收益必须按完整任务核算}


\subsection{三条降本提速路径}

本项目不以减少必要的业务处理换取成本下降，而是将“需要生成内容”与“只需作出判断”的任务分开。以下收益均以业务质量达标为前提；原有调用方式和费用应先通过实际日志确认，不将示例流程直接视为现网事实。


{\small\setlength{\tabcolsep}{4pt}\renewcommand{\arraystretch}{1.28}
\begin{longtable}{>{\raggedright\arraybackslash}p{0.150\dimexpr\linewidth-8\tabcolsep\relax}>{\raggedright\arraybackslash}p{0.240\dimexpr\linewidth-8\tabcolsep\relax}>{\raggedright\arraybackslash}p{0.290\dimexpr\linewidth-8\tabcolsep\relax}>{\raggedright\arraybackslash}p{0.320\dimexpr\linewidth-8\tabcolsep\relax}}
\caption{三条降本提速路径及生效条件}\label{tab:saving-paths}\\
\toprule
\textbf{优化路径} & \textbf{原有资源消耗} & \textbf{本方案如何减少消耗} & \textbf{提速来源与生效条件} \\
\midrule
\endfirsthead
\toprule
\textbf{优化路径} & \textbf{原有资源消耗} & \textbf{本方案如何减少消耗} & \textbf{提速来源与生效条件} \\
\midrule
\endhead
\midrule\multicolumn{4}{r}{\scriptsize\color{Muted}续下页}\\\endfoot
\bottomrule\endlastfoot
替换昂贵判断 & 为获取一个类别或是非结果，调用通用生成式大模型 & 对适合的节点采用分类或选项打分后端，直接返回结构化决策，替代原生成式判断调用 & 减少生成式解码及相关处理；只有候选后端的实际推理、部署和回退开销更低时才有净收益 \\
减少生成调用 & 默认改写流程对无需改写的问题也执行一次生成 & 先判断是否需要改写，明确完整的问题跳过改写生成，必要改写仍保留 & 无需改写的请求少等待一次生成；节省必须覆盖新增判断开销，且不能损害必要改写的召回 \\
缩短生成输入 & 冗余检索片段进入回答模型，增加输入 token & 筛选与问题有关且能支撑回答的证据，保留必要的例外和冲突材料，减少冗余上下文 & 降低输入处理量及相应费用；筛选本身也消耗资源，实际净时延受后端、批处理和缓存影响 \\
\end{longtable}}

三条路径可以独立采用，不要求同时替换所有节点。输出更短不等于模型一定更小或更快；若所选后端仍需生成文本，则应按实际生成过程核算，不将其视为无生成开销的分类器。


\subsection{同一业务问题的使用前后对比}

以“退款已经申请了，怎么还没到账？”为例，演示一条包含意图识别、问题改写、规则检索和回答生成的链路。假设原流程分别调用生成模型完成意图识别、默认改写和最终回答，并将检索返回的 5 个片段全部送入回答模型。实际业务查询等其他必要步骤保持不变。


\textbf{以下数字仅用于说明链路变化，不是实测结果，也不是对现网调用次数的确认。}


{\small\setlength{\tabcolsep}{4pt}\renewcommand{\arraystretch}{1.28}
\begin{longtable}{>{\raggedright\arraybackslash}p{0.220\dimexpr\linewidth-6\tabcolsep\relax}>{\raggedright\arraybackslash}p{0.380\dimexpr\linewidth-6\tabcolsep\relax}>{\raggedright\arraybackslash}p{0.400\dimexpr\linewidth-6\tabcolsep\relax}}
\caption{示例业务链路优化前后对照}\label{tab:request-example}\\
\toprule
\textbf{环节} & \textbf{使用前：示例基线} & \textbf{使用后：预期路径} \\
\midrule
\endfirsthead
\toprule
\textbf{环节} & \textbf{使用前：示例基线} & \textbf{使用后：预期路径} \\
\midrule
\endhead
\midrule\multicolumn{3}{r}{\scriptsize\color{Muted}续下页}\\\endfoot
\bottomrule\endlastfoot
意图识别 & 1 次生成式调用，输出退款进度查询类别 & 决策模型判断；满足接受策略后进入同一业务流程，不再执行原生成式分类调用 \\
问题改写 & 1 次生成式调用，对问题进行改写 & 决策模型确认问题已完整明确后跳过改写；需要改写时仍调用生成模型 \\
检索上下文 & 5 个检索片段全部进入回答模型 & 假设经评估其中 3 个已覆盖必要证据，则保留这 3 个；证据不足时不得机械裁剪 \\
回答生成 & 1 次生成式调用，读取全部候选片段 & 仍保留 1 次回答生成调用，读取筛选后的证据 \\
调用结构 & 共 3 次生成式调用 & 在判断可靠且无需改写的示例路径中，变为 1 次生成式调用，另增加决策与筛选开销 \\
\end{longtable}}

这一示例体现的是：\textbf{保留必要的回答生成，把用于判断的生成调用替换掉，把不必要的改写调用省下来，把冗余证据输入减下来。} 片段筛选的调用次数取决于后端与批处理方式，不假定多个片段可零开销或一次完成。


生成式调用从 3 次变为 1 次，不等于总费用或完整链路时延必然下降三分之二。问题需要改写、发生回退或筛选开销过大时，收益会缩小，甚至不成立；最终按完整任务的实际费用与时延判断。


\subsection{双闭环 RSI 如何支持持续降本}

仅替换模型可能因业务表达识别不准而增加误判、澄清和回退，抵消轻量化收益。双闭环 RSI 的作用，是在固定业务含义和质量约束下，持续改善低成本模型承接判断任务的能力，而不是在线上增加一轮自优化操作。


\begin{itemize}

\item \textbf{提示优化闭环}：针对“退款到账”与“结算到账”等混淆，优化类别描述、排除说明和示例，争取在不训练模型的情况下改善判断质量。

\item \textbf{业务微调闭环}：当描述已清晰、输入与标签已修复，仍存在稳定的领域误判时，通过审核后的业务数据补强识别能力，并与未微调版本对照。

\end{itemize}

只有在接受部分错误率满足约束的前提下，提高自动接受覆盖率、减少不必要的回退才算有效改善，不能靠降低接受标准获得表面上的节省。两条闭环都在线下运行，候选通过评测和审批后才发布；标注、搜索、训练及验收费用按约定周期摊销，不能从收益核算中删除。


\subsection{预期成效与参赛验证成果}

项目预期实现“业务质量达标、完整任务成本下降、完整链路时延缩短”，并分别验证直接替换、提示优化和业务微调的贡献。当前文档未提供实测数据，以下为待验证目标，不作为已实现成果。


{\small\setlength{\tabcolsep}{4pt}\renewcommand{\arraystretch}{1.28}
\begin{longtable}{>{\raggedright\arraybackslash}p{0.150\dimexpr\linewidth-8\tabcolsep\relax}>{\raggedright\arraybackslash}p{0.240\dimexpr\linewidth-8\tabcolsep\relax}>{\raggedright\arraybackslash}p{0.290\dimexpr\linewidth-8\tabcolsep\relax}>{\raggedright\arraybackslash}p{0.320\dimexpr\linewidth-8\tabcolsep\relax}}
\caption{项目预期成效与验收关注点}\label{tab:expected-outcomes}\\
\toprule
\textbf{维度} & \textbf{使用前需记录的基线} & \textbf{使用后预期成效} & \textbf{验收关注点} \\
\midrule
\endfirsthead
\toprule
\textbf{维度} & \textbf{使用前需记录的基线} & \textbf{使用后预期成效} & \textbf{验收关注点} \\
\midrule
\endhead
\midrule\multicolumn{4}{r}{\scriptsize\color{Muted}续下页}\\\endfoot
\bottomrule\endlastfoot
业务质量 & 当前路由正确率、关键类别召回和最终回答质量 & 在预先约定的允许退化范围内达标，关键业务指标不突破下限 & 不以漏改写、误删证据或强行分类换取节省 \\
完整任务成本 & 当前每成功任务总成本及费用构成 & 将新增决策、回退、部署运维与离线优化摊销一并计入后，成本仍下降 & 区分节点费用下降与完整任务净节省 \\
响应效率 & 相同负载下的完整链路 p50/p95、吞吐和超时率 & 完整链路 p95 缩短，超时率不恶化，并报告其他性能指标 & 不能只用某个决策节点的加速代表用户整体体验 \\
优化可持续性 & 当前模型、提示、数据与人工调整记录 & 形成可复现的提示优化与可选微调实验，发布可追溯、可回退的版本 & 无有效增量时保留原版本，不强行上线候选 \\
\end{longtable}}

成本降幅按“（基线每成功任务总成本－优化后每成功任务总成本）÷ 基线每成功任务总成本”计算；时延降幅按相同口径的完整链路 p95 比较。两者分别报告，不能用生成调用次数或输入 token 的降幅替代。具体质量门槛和验收目标在查看候选结果前确定，详细数据划分、测试条件及结果表见第六章。


参赛验证拟形成三类成果：可运行的首个意图识别试点、可复现的双闭环优化与对照实验、包含质量和完整成本时延的前后对比报告。扩展场景按实际验证进度分别呈现；未完成部分明确标注为计划或待验证，不与已完成成果混写。


\section{总体架构}

系统分为线上决策和离线优化两部分。线上只加载经过批准的版本，离线负责分析错误、生成候选并验证改进；两者通过版本包连接。



\begin{center}\resizebox{\linewidth}{!}{\begin{tikzpicture}[x=1mm,y=1mm]
\node[dark,text width=145mm] at (59,22) {线上决策：只加载经过批准的版本};
\node[box,text width=28mm] (b) at (0,0) {业务节点};
\node[accent,text width=28mm] (s) at (38,0) {决策服务\\模型适配与推理};
\node[box,text width=32mm] (p) at (78,0) {接受／澄清／回退\\业务路由或原流程};
\node[box,text width=28mm] (l) at (118,0) {调用记录\\质量反馈};
\draw[arr] (b.east)--(s.west);\draw[arr](s.east)--(p.west);\draw[arr](p.east)--(l.west);
\node[dark,text width=145mm] at (59,-28) {离线优化：审核错误、生成候选、独立验证};
\node[box,text width=28mm] (d) at (0,-50) {反馈审核\\错误分析};
\node[accent,text width=28mm] (o) at (38,-50) {提示优化\\业务微调};
\node[box,text width=32mm] (e) at (78,-50) {统一评测\\独立验证与校准\\冻结候选并盲测};
\node[box,text width=28mm] (v) at (118,-50) {人工审批\\不可变版本包};
\draw[arr](d.east)--(o.west);\draw[arr](o.east)--(e.west);\draw[arr](e.east)--(v.west);
\draw[arr,draw=Muted,line width=.8pt] (l.south)--(118,-13)--(142,-13)--(142,-39)--node[lab,above,pos=.45]{业务反馈进入审核}(0,-39)--(d.north);
\draw[arr,draw=Teal,line width=.9pt] (v.south)--(118,-65)--(-24,-65)--(-24,11)--node[lab,above,pos=.6]{批准版本加载}(38,11)--(s.north);
\end{tikzpicture}}\end{center}
\figcap{总体架构：线上执行与离线优化分离，版本包连接两侧}


核心组件及其职责如下：


{\small\setlength{\tabcolsep}{4pt}\renewcommand{\arraystretch}{1.28}
\begin{longtable}{>{\raggedright\arraybackslash}p{0.420\dimexpr\linewidth-4\tabcolsep\relax}>{\raggedright\arraybackslash}p{0.580\dimexpr\linewidth-4\tabcolsep\relax}}
\caption{系统组件与职责分工}\label{tab:system-components}\\
\toprule
\textbf{组件} & \textbf{主要职责} \\
\midrule
\endfirsthead
\toprule
\textbf{组件} & \textbf{主要职责} \\
\midrule
\endhead
\midrule\multicolumn{2}{r}{\scriptsize\color{Muted}续下页}\\\endfoot
\bottomrule\endlastfoot
任务与数据管理 & 保存业务类别、输入要求、标注与数据划分，固定评测依据 \\
决策服务与模型适配 & 接入候选模型，统一结构化结果，处理超时与后端差异 \\
提示优化器 & 分析允许查看的错误，提出并比较提示候选 \\
微调执行器 & 基于审核数据与固定模板运行训练实验 \\
评测与版本管理 & 计算质量、成本和时延，记录血缘，管理发布与回退 \\
\end{longtable}}

这些组件可以先通过配置文件、批处理脚本、离线评测工具和轻量服务实现，不要求建设完整训练平台。候选模型可从 Kev、Laya、D3-Flash 等选择，中文效果、描述输入能力、训练支持和许可需逐一核实；规则与轻量分类器也应纳入对照。


线上调用只需提供任务标识和业务输入；类别描述、模型、阈值等由服务端版本包确定，避免调用方随意改变任务含义。


\section{核心实现方案}

\subsection{业务节点如何替换}

以意图识别为例，调用过程分为四步：


\begin{enumerate}

\item \textbf{准备输入}：提取当前问题和决定意图所需的前文，保留必要产品、渠道等授权信息。

\item \textbf{执行判断}：模型适配层将任务定义转为后端支持的格式，获取类别及原始分数。

\item \textbf{应用策略}：根据输入完整性、支持范围及校准后的接受条件，返回接受、澄清、回退或错误状态。

\item \textbf{记录结果}：关联请求、任务版本、模型/提示版本、输出、耗时、费用与回退原因，供评测和追踪。

\end{enumerate}

不同模型的原始分数含义可能不同，应分别保留并校准，不能直接混用。模型给出较高分数，也不能保证输入属于支持范围或业务结果一定正确。


对“什么时候到？”这样的信息缺失，优先澄清；对超时、输入超限或不可靠结果，按策略回到原流程。决策与回退共用总截止时间，避免增加不可控等待。回退大模型的输出同样需要评测。


\subsection{提示优化闭环如何运行}

提示优化固定模型权重和业务类别含义，只调整判断指令、类别描述、包含/排除说明及允许的示例。


例如，初始类别描述为：


\begin{insight}{类别描述示例}退款到账问题。\end{insight}

如果训练池中反复出现退款与商户结算混淆，优化器可以提出：


\begin{insight}{类别描述示例}查询已发起的消费者退款何时到账或当前进度；不包含商户收款结算，也不包含尚未发起的退款申请。\end{insight}

实现时，优化器读取经审核的训练错误和混淆统计，提出少量候选，并说明改动与预期解决的问题。执行器检查类别集合、字段、长度和示例来源，再在固定开发集上比较初始版本、当前最优与新候选。


只有满足质量约束且产生有效收益的候选才被保留；连续无改善或达到调用费用、轮次、时间上限即停止。优化器不能修改类别定义、标签或评测程序，关键语义变化需人工审查。


该闭环仅适用于支持指令或类别描述输入的后端。固定标签分类头不能因为修改外部描述就自动获得能力提升。


\subsection{业务微调闭环如何运行}

微调解决的是模型对业务表达的稳定识别能力不足，不能代替输入修复和标签治理。


{\small\setlength{\tabcolsep}{4pt}\renewcommand{\arraystretch}{1.28}
\begin{longtable}{>{\raggedright\arraybackslash}p{0.420\dimexpr\linewidth-4\tabcolsep\relax}>{\raggedright\arraybackslash}p{0.580\dimexpr\linewidth-4\tabcolsep\relax}}
\caption{典型错误现象与优先处理方式}\label{tab:error-diagnosis}\\
\toprule
\textbf{错误现象} & \textbf{优先处理} \\
\midrule
\endfirsthead
\toprule
\textbf{错误现象} & \textbf{优先处理} \\
\midrule
\endhead
\midrule\multicolumn{2}{r}{\scriptsize\color{Muted}续下页}\\\endfoot
\bottomrule\endlastfoot
缺少前文，无法确定“到账”指什么 & 补充上下文或澄清 \\
同类样本标注不一致、类别重叠 & 修正规范并仲裁标签 \\
类别基本正确，但接受与回退不合理 & 校准分数和调整策略 \\
描述已明确，仍稳定误判业务术语或口语 & 补充真实标注，比较微调 \\
输入被截断或超过模型长度限制 & 调整输入方案或更换后端 \\
\end{longtable}}

进入微调后，先冻结任务与提示，按错误类型提出补数需求。新数据完成授权、脱敏、标注与去重后，使用批准的训练模板生成少量配置，在限定 GPU 时间和候选数量内执行。


训练方式以所选模型支持为准，可比较决策头、适配器或其他成熟微调方式，不假设所有后端都支持同一种 LoRA。除困难样本外，还需保留日常分布数据，防止局部提升伴随通用问题退化。


训练完成后，与未微调版本使用同样数据、输入和提示比较，并检查不同随机种子的稳定性。微调无收益则保留原版本，不强推训练产物。


\subsection{两条闭环如何协同}

初期先固定模型优化提示，再固定任务与入围提示比较微调。后续可以交替运行，但每次应明确主要改动，避免同时修改提示、训练数据和阈值，无法解释收益来源。



\begin{center}\begin{tikzpicture}
\node[dark,text width=53mm] (f) {经审核的真实错误 → 错误诊断};
\node[box,below left=10mm and 8mm of f,text width=37mm] (fix) {输入或标签问题\\修复输入与标注};
\node[accent,below=10mm of f,text width=37mm] (p) {描述表达不清\\固定模型，优化提示};
\node[accent,below right=10mm and 8mm of f,text width=37mm] (t) {领域能力不足\\冻结任务与提示，补数微调};
\draw[arr] (f)--(fix);\draw[arr](f)--(p);\draw[arr](f)--(t);
\node[box,below=12mm of p,text width=55mm] (e) {固定数据与预算下评测\\满足质量约束且有收益？};
\draw[arr](p)--(e);\draw[arr](t)|-(e);
\node[box,left=7mm of e,text width=35mm] (k) {保留原版本\\或停止实验};
\draw[arr](e)--node[lab,above]{否}(k);
\node[accent,below=9mm of e,text width=55mm] (v) {独立验证与策略校准\\冻结候选并盲测};
\node[dark,below=7mm of v,text width=55mm] (a) {人工审批 → 发布 → 收集反馈};
\draw[arr](e)--node[lab,right]{是}(v);\draw[arr](v)--(a);
\draw[arr] (a.east)--++(35mm,0)|-node[lab,near start,right]{审核后进入下一轮}(f.east);
\end{tikzpicture}\end{center}
\figcap{双闭环协同：按错误成因选择优化路径，验证收益后发布}


最后一条反馈边表示发布后继续采集业务反馈，须经审核才能进入下一轮；不表示模型会在线自动训练或直接改写生产配置。


\subsection{评测、版本与运行保障}

任务规范、数据、模型、提示、输入处理、校准与接受策略共同组成一个版本包。实验记录父版本、改进假设、配置、指标及费用；上线时整包切换，异常时整包回退。


最终候选通过离线验收后，先做不改变业务结果的影子验证，再经批准灰度。监控关键类别错误、接受覆盖率、回退率、超时、成本和输入分布变化；质量或运行指标触发预设条件时暂停新版本。


数据、训练、优化器和评测分别授权；训练只运行批准模板。业务语义变更与生产发布保留人审，决策结果不直接授权退款、资金操作或风控处置。


\section{效果验证方案}

\subsection{验证要回答的问题}


\begin{center}\begin{tikzpicture}
\node[box,text width=24mm] (a) {A\quad 现有模型\\\footnotesize 质量与费用基线};
\node[accent,right=5mm of a,text width=24mm] (b) {B\quad 原始决策\\\footnotesize 直接替换收益};
\node[accent,right=5mm of b,text width=24mm] (c) {C\quad 提示优化\\\footnotesize 描述改进增量};
\node[box,right=5mm of c,text width=24mm] (d) {D\quad 可选微调\\\footnotesize 领域能力增量};
\node[dark,right=5mm of d,text width=24mm] (e) {E\quad 最终策略\\\footnotesize 完整链路收益};
\draw[arr](a)--(b);\draw[arr](b)--(c);\draw[arr](c)--(d);\draw[arr](d)--(e);
\end{tikzpicture}\end{center}
\figcap{实验递进关系；E 使用开发与验证阶段选定的候选，D 可按条件开展}


验证围绕三件事展开：基础替换是否值得，双闭环是否带来增量，加入回退后的完整业务链路是否仍有收益。


{\small\setlength{\tabcolsep}{4pt}\renewcommand{\arraystretch}{1.28}
\begin{longtable}{>{\raggedright\arraybackslash}p{0.220\dimexpr\linewidth-6\tabcolsep\relax}>{\raggedright\arraybackslash}p{0.380\dimexpr\linewidth-6\tabcolsep\relax}>{\raggedright\arraybackslash}p{0.400\dimexpr\linewidth-6\tabcolsep\relax}}
\caption{决策模型优化对照组设置}\label{tab:evaluation-groups}\\
\toprule
\textbf{对照组} & \textbf{方案} & \textbf{主要验证内容} \\
\midrule
\endfirsthead
\toprule
\textbf{对照组} & \textbf{方案} & \textbf{主要验证内容} \\
\midrule
\endhead
\midrule\multicolumn{3}{r}{\scriptsize\color{Muted}续下页}\\\endfoot
\bottomrule\endlastfoot
A & 当前生成式大模型与现有提示 & 现有质量、费用和时延 \\
B & 原始决策模型与初始描述 & 直接替换收益 \\
C & B + 提示优化 & 提示优化的增量 \\
D & 优化提示 + 业务微调 & 在 C 基础上是否进一步改善 \\
E & 入围候选 + 接受/澄清/回退策略 & 最终业务链路表现 \\
\end{longtable}}

A/B/C 为基础对照；D 在具备有效数据和训练条件时开展。若要单独归因微调贡献，可增加“初始描述 + 微调”消融组；也可让现有大模型获得同等提示优化预算，检验收益是否主要来自描述改进。


E 应评估经过开发/验证阶段选定的候选，不能在盲测中反复挑选冠军。所有方案使用相同业务语义和必要上下文。


\subsection{测试数据如何准备}

以真实、人工标注的业务问题为主，覆盖正常表达、易混淆类别、多轮指代、口语与错别字、未知意图、信息不足和多意图。类别不均衡时，同时保留真实分布集和困难集，分别报告。


训练/错误分析、搜索开发、独立验证、校准和最终盲测按用途隔离；同一会话、事件和近重复样本不能跨集合。优化器只查看允许的数据，盲测不用于生成候选或调阈值。


数据规模根据类别覆盖与要证明的差异确定，不预先用固定条数保证效果。可先用小批样本完善标注边界，再扩展独立评测；合成数据可辅助训练，不充当真实业务验收数据。


\subsection{离线同题对照}

对各组重放相同问题，记录预测结果、耗时和费用，并对错误按类别和场景拆分。


{\small\setlength{\tabcolsep}{4pt}\renewcommand{\arraystretch}{1.28}
\begin{longtable}{>{\raggedright\arraybackslash}p{0.220\dimexpr\linewidth-6\tabcolsep\relax}>{\raggedright\arraybackslash}p{0.380\dimexpr\linewidth-6\tabcolsep\relax}>{\raggedright\arraybackslash}p{0.400\dimexpr\linewidth-6\tabcolsep\relax}}
\caption{离线评测指标及其用途}\label{tab:offline-metrics}\\
\toprule
\textbf{维度} & \textbf{核心指标} & \textbf{目的} \\
\midrule
\endfirsthead
\toprule
\textbf{维度} & \textbf{核心指标} & \textbf{目的} \\
\midrule
\endhead
\midrule\multicolumn{3}{r}{\scriptsize\color{Muted}续下页}\\\endfoot
\bottomrule\endlastfoot
分类质量 & Accuracy、Macro-F1、各类召回/精确率 & 避免整体分数掩盖关键类别退化 \\
接受策略 & 自动接受覆盖率、接受部分错误率 & 同时衡量“处理多少”和“处理得多可靠” \\
性能 & p50/p95、吞吐、超时率 & 判断实际并发下是否提速 \\
费用 & 单节点费用、完整任务费用 & 区分节点节省与总体节省 \\
\end{longtable}}

使用真实中文长度、类别数、设备和并发，区分冷启动与稳定运行；同时纳入现有模型的缓存和批处理条件。对质量差值报告合适的区间，按会话分组进行配对分析；较小提升需足够样本支撑，训练与推理随机性分别检查。


未知输入和关键业务错误单独观察，不能只筛选高置信度成功样本展示准确率。


\subsection{完整业务链路验证}

节点评测通过后，还要验证下游结果，防止“分类更便宜，业务反而变差”。


{\small\setlength{\tabcolsep}{4pt}\renewcommand{\arraystretch}{1.28}
\begin{longtable}{>{\raggedright\arraybackslash}p{0.420\dimexpr\linewidth-4\tabcolsep\relax}>{\raggedright\arraybackslash}p{0.580\dimexpr\linewidth-4\tabcolsep\relax}}
\caption{完整业务链路场景与下游验证}\label{tab:chain-validation}\\
\toprule
\textbf{场景} & \textbf{下游验证方式} \\
\midrule
\endfirsthead
\toprule
\textbf{场景} & \textbf{下游验证方式} \\
\midrule
\endhead
\midrule\multicolumn{2}{r}{\scriptsize\color{Muted}续下页}\\\endfoot
\bottomrule\endlastfoot
意图识别 & 检查最终路由正确率、误转、用户纠正和转人工情况 \\
是否改写 & 比较原始与改写后检索的关键证据召回、回答质量及完整任务费用 \\
片段筛选 & 检查关键证据是否误删、回答是否有依据和完整，并与现有 reranker 比较 \\
\end{longtable}}

先通过影子流量做同题比较，再按会话固定分组进行受控灰度。影子双跑的额外开销单列为实验费用。线上没有投诉不等于正确，应结合抽样审核及可信反馈评估。


\subsection{成本口径与验收结论}

成本至少包括决策推理、原流程回退、仍需执行的生成/检索，以及新增部署运维；标注、提示搜索、训练和验收费用按约定周期摊销。


\begin{insight}{完整成本口径}\[\text{每成功任务总成本}=\frac{\text{全部任务总成本}}{\text{成功任务数}}\]\end{insight}

回退和失败任务的支出不能排除。分别报告试点规模与预期业务量下的成本；低流量时，自部署固定费用可能抵消推理节省。


建议最终用一张结果表呈现，实测前统一标记“待测”：


{\small\setlength{\tabcolsep}{4pt}\renewcommand{\arraystretch}{1.28}
\begin{longtable}{>{\raggedright\arraybackslash}p{0.280\dimexpr\linewidth-12\tabcolsep\relax}>{\raggedright\arraybackslash}p{0.144\dimexpr\linewidth-12\tabcolsep\relax}>{\raggedright\arraybackslash}p{0.144\dimexpr\linewidth-12\tabcolsep\relax}>{\raggedright\arraybackslash}p{0.144\dimexpr\linewidth-12\tabcolsep\relax}>{\raggedright\arraybackslash}p{0.144\dimexpr\linewidth-12\tabcolsep\relax}>{\raggedright\arraybackslash}p{0.144\dimexpr\linewidth-12\tabcolsep\relax}}
\caption{方案对照结果记录表（待测）}\label{tab:results-template}\\
\toprule
\textbf{指标} & \textbf{当前大模型} & \textbf{原始决策模型} & \textbf{提示优化版} & \textbf{微调版（可选）} & \textbf{最终策略} \\
\midrule
\endfirsthead
\toprule
\textbf{指标} & \textbf{当前大模型} & \textbf{原始决策模型} & \textbf{提示优化版} & \textbf{微调版（可选）} & \textbf{最终策略} \\
\midrule
\endhead
\midrule\multicolumn{6}{r}{\scriptsize\color{Muted}续下页}\\\endfoot
\bottomrule\endlastfoot
分类质量与区间 & 待测 & 待测 & 待测 & 待测 & 待测 \\
关键类别召回 & 待测 & 待测 & 待测 & 待测 & 待测 \\
自动接受覆盖率 / 错误率 & 按现有策略记录 & 待测 & 待测 & 待测 & 待测 \\
每任务生成式调用次数（含回退） & 待测 & 待测 & 待测 & 待测 & 待测 \\
回答模型每次调用输入 token（均值 / p95） & 待测 & 待测 & 待测 & 待测 & 待测 \\
完整链路 p95 & 待测 & 待测 & 待测 & 待测 & 待测 \\
每成功任务总成本 & 待测 & 待测 & 待测 & 待测 & 待测 \\
回退率与费用 & 按现有策略记录 & 待测 & 待测 & 待测 & 待测 \\
\end{longtable}}

业务质量的允许退化范围、关键类别召回下限、覆盖率要求及 SLA 在看候选结果前确定。\textbf{先满足质量约束，再比较成本和时延。} 降本提速比例以实测为准，样本不足时仅报告试验观察，不宣称已证明生产收益。


\section{预期价值与适用边界}

本方案的价值有三个层面：通过适配判断任务的模型减少不必要的生成开销；用真实错误推动可验证的持续改进；复用任务定义、模型适配和评测能力，降低后续节点接入成本。


适合优先尝试的节点应具有明确输出、可获得可靠标签、稳定业务边界和足够调用量。开放式生成、复杂推理或缺乏必要上下文的任务，不宜仅靠分类替换。


项目创新重点不是单独使用提示优化或微调，而是把支付业务边界、双闭环优化、质量约束和完整成本核算连起来，形成从真实错误到可追溯发布版本的链路。


如果实测发现回退率过高、质量不达标或总成本无改善，应只替换有效节点、保留原模型，或停止相应优化。有效结论也可以是“提示优化有用，微调暂时无收益”。


\textbf{候选能力核查入口：} \href{https://github.com/jaredpalmer/kev}{Kev}、\href{https://github.com/NandhaKishorM/laya/blob/main/docs/finetune.md}{Laya 微调}、\href{https://huggingface.co/vllm-sr/d3-flash}{D3-Flash}。链接沿用原稿，实施前重新核实版本、训练方式和许可。


\end{document}
